% ID: SEL-D0015-BB-E03
% ESTADO: cotejada. Segunda pasada visual completa y borrador compilado revisados.
% Fuente: Selectividad/matematicas_aplicadas/matematicas_aplicadas_P_C.pdf
% SHA-256: 463bf63d4697a18a0b9efc9287e7b0a593c803bab98a3cfab0e260067a0dd02c
% Página: 1; bloque B; ejercicio original 3.

Se considera la función $f(x)=ax^3+bx+4$, con $a$ y $b$ números reales.
\begin{enumerate}[label=\alph*)]
  \item \textbf{(1 puntos)} Determine los valores $a$ y $b$ para que $f$
  tenga un extremo relativo en el punto $(2,36)$.
  \item \textbf{(0.75 puntos)} Para $a=4$ y $b=-3$, estudie la monotonía de
  $f$ y determine sus extremos relativos.
  \item \textbf{(0.75 puntos)} Para $a=4$ y $b=-3$, calcule la función $F(x)$
  que verifica $F'(x)=f(x)$ y $F(2)=10$.
\end{enumerate}
